\documentclass[11pt]{amsart}
\usepackage{amsmath,amssymb,amsthm}
\usepackage[margin=1.15in]{geometry}
\usepackage[hidelinks]{hyperref}
\newtheorem{theorem}{Theorem}[section]
\newtheorem{corollary}[theorem]{Corollary}
\newtheorem{lemma}[theorem]{Lemma}
\theoremstyle{remark}
\newtheorem{remark}[theorem]{Remark}
\newcommand{\doi}[1]{\href{https://doi.org/#1}{doi:\nolinkurl{#1}}}
\newcommand{\arxiv}[1]{\href{https://arxiv.org/abs/#1}{arXiv:\nolinkurl{#1}}}
\newcommand{\gap}{\operatorname{gap}}
\setlength{\emergencystretch}{2.5em}

\begin{document}

\title[Single-card gap of the overlapping-cycles shuffle]{An upper bound for the single-card spectral gap of the
overlapping-cycles shuffle}
\author{Yitzchak Shmalo}
\address{Einstein Institute of Mathematics, The Hebrew University of Jerusalem, Jerusalem, Israel}
\email{yspennstate@gmail.com}
\subjclass[2020]{60J10, 11H06, 11J70}
\keywords{card shuffling, spectral gap, overlapping-cycles shuffle, Hermite constant, Diophantine approximation}

\begin{abstract}
In the overlapping-cycles shuffle of a deck of $n$ cards, the card at position $n$ or the card at position $n-k$ is
moved to the top, each with probability one half. Angel, Peres and Wilson determined the spectral gap $\gap(n,k)$ of
the chain followed by a single card, up to a factor $1+o(1)$ uniform in $k$, as the minimum of an explicit family of
quadratic expressions. They asked whether $\gap(n,k)\le(2\pi^2/\sqrt3+o(1))\,k^{1/2}n^{-2}$ for all $1\le k<n$. We show
that it does. Their minimum is a constant multiple of the squared length of the shortest vector of a plane lattice of
determinant $2$, and Hermite's bound for plane lattices gives the inequality with this constant; Dirichlet's
approximation theorem, used by Angel, Peres and Wilson, gives $2\pi^2$. At $n=26$, $k=12$ the lattice is hexagonal and the bound on the
minimum is attained exactly.
\end{abstract}

\maketitle

\section{Introduction}\label{sec:intro}

The overlapping-cycles shuffle with parameters $n$ and $1\le k<n$ acts on a deck of $n$ cards: at each step, with
probability $1/2$ each, the card at position $n$ or the card at position $n-k$ is moved to the top of the deck. Jonasson
first studied it~\cite{Jon06}, and Angel, Peres and Wilson~\cite{APW08} determined the
spectral gap $\gap(n,k)$ of the Markov chain on $\{1,\dots,n\}$ followed by the position of a single card, that is,
$1$ minus the largest modulus of an eigenvalue other than $1$ of its transition matrix (the chain is not reversible,
and this is the convention of~\cite{APW08}). The answer
depends on the Diophantine properties of $k/n$: the relaxation time is of order $n^2$ when $k/n$ is near a rational
with small denominator and of order $n^{3/2}$ when $k/n$ is badly approximable. Mixing of the whole deck was studied
by Jonasson~\cite{Jon11} and, more recently, by Blumberg, Morris and Oberschelp~\cite{BMO23}.

The main theorem of~\cite{APW08} expresses the gap through the quantity
\begin{equation}\label{eq:gamma}
\gamma(n,k,m)=\frac{\pi^2}{2n^2}\Bigl(\frac{m^2k}{n}+4n\Bigl\|\frac{m(n-k)}{2n}\Bigr\|^2\Bigr),\qquad
\gamma(n,k)=\min_{m\in\mathbb Z,\,m\ne0}\gamma(n,k,m),
\end{equation}
where $\|x\|=\operatorname{dist}(x,\mathbb Z)$.

\begin{theorem}[Angel, Peres and Wilson~{\cite[Theorem~4]{APW08}}]\label{thm:apw}
For all large enough $n$ and every $1\le k\le n-1$,
\[
\gap(n,k)=(1+o(1))\,\gamma(n,k),
\]
where the $o(1)$ tends to zero uniformly in $k$ as $n\to\infty$.
\end{theorem}

Their proof of Theorem~\ref{thm:apw}~\cite[p.~13]{APW08} applies Dirichlet's approximation theorem to
$(1-k/n)/2$ and gives
\[
\gamma(n,k)\le 2\pi^2\,\frac{\sqrt k}{n^2},
\]
and with it $\gap(n,k)\le(2\pi^2+o(1))n^{-3/2}$ for every $k$~\cite[Theorem~2(i)]{APW08}. They
remarked that this bound should improve by a factor of $\sqrt3$, noted that $\gamma(n,k)=(2\pi^2/\sqrt3+o(1))n^{-3/2}$
for some $k=n-(2/3^{1/4}+o(1))n^{1/4}$, so that such an improvement would be sharp for $\gamma$, and asked in their second open problem whether
$\gap(n,k)\le(2\pi^2/\sqrt3+o(1))k^{1/2}n^{-2}$ holds. They observed that $\gamma(n,k)$ appears to be no larger than
$(2\pi^2/\sqrt3+o(1))k^{1/2}n^{-2}$, as their Figure~5 illustrates for $n=1000$~\cite[Section~6]{APW08}. The answer
is yes, and the observation holds as an exact inequality for $\gamma$.

\begin{theorem}\label{thm:main}
For every $n\ge2$ and every $1\le k\le n-1$,
\begin{equation}\label{eq:bound}
\gamma(n,k)\le\frac{2\pi^2}{\sqrt3}\,\frac{\sqrt k}{n^2}.
\end{equation}
Consequently $\gap(n,k)\le(2\pi^2/\sqrt3+o(1))\,k^{1/2}n^{-2}$, with $o(1)\to0$ uniformly in $1\le k<n$.
\end{theorem}

The inequality~\eqref{eq:bound} is exact and holds for every $n$; only the statement about $\gap(n,k)$ uses
Theorem~\ref{thm:apw}. Since $k<n$, Theorem~\ref{thm:main} also gives $\gap(n,k)\le(2\pi^2/\sqrt3+o(1))n^{-3/2}$, which is Theorem~2(i)
of~\cite{APW08} improved by the factor $\sqrt3$. The proof replaces Dirichlet's theorem by one observation: $\gamma(n,k)$ is a
multiple of the minimum of a positive definite binary quadratic form of fixed determinant, so the bound of Lagrange
and Gauss for binary forms (Lemma~\ref{lem:hermite}) applies. Section~\ref{sec:checks} checks~\eqref{eq:bound} for $n\le2000$ and compares
$\gamma(n,k)$ with the exact gap; both computations are made with the program of~\cite{Repo}.

\section{Proof}\label{sec:proof}

Write $\theta=1-k/n$ and $A=\sqrt k/n$, so that $0<A<1$. For $m\in\mathbb Z$ let
$b_m=\operatorname{dist}(m\theta,2\mathbb Z)$. Since $\operatorname{dist}(x/2,\mathbb Z)=\tfrac12\operatorname{dist}(x,2\mathbb Z)$,
we have $4n\|m(n-k)/(2n)\|^2=4n\|m\theta/2\|^2=n\,b_m^2$, and therefore
\begin{equation}\label{eq:rewrite}
\gamma(n,k,m)=\frac{\pi^2}{2n^2}\Bigl(\frac{m^2k}{n}+n\,b_m^2\Bigr)
=\frac{\pi^2\sqrt k}{2n^2}\Bigl(A\,m^2+\frac{b_m^2}{A}\Bigr),
\end{equation}
because $\sqrt k\,(Am^2+b_m^2/A)=m^2k/n+n\,b_m^2$.

Let $\Lambda\subset\mathbb R^2$ be the lattice spanned by
\[
u_1=\Bigl(\sqrt A,\ \frac{\theta}{\sqrt A}\Bigr),\qquad u_2=\Bigl(0,\ \frac{2}{\sqrt A}\Bigr).
\]
Its determinant is $|\det(u_1,u_2)|=2$. The vector $mu_1-ju_2$ has squared length
$Am^2+(m\theta-2j)^2/A$, and for fixed $m$ the minimum over $j\in\mathbb Z$ is $Am^2+b_m^2/A$. By~\eqref{eq:rewrite},
\begin{equation}\label{eq:lattice}
\gamma(n,k)=\frac{\pi^2\sqrt k}{2n^2}\,\mu,\qquad
\mu=\min\bigl\{|mu_1-ju_2|^2:\ m,j\in\mathbb Z,\ m\ne0\bigr\},
\end{equation}
so $\mu$ is the smallest squared length of a vector of $\Lambda$ with nonzero first coordinate.

\begin{lemma}[Lagrange, Gauss]\label{lem:hermite}
Every lattice in $\mathbb R^2$ of determinant $d$ contains a nonzero vector $v$ with $|v|^2\le(2/\sqrt3)\,d$, with
equality for the shortest vector exactly when the lattice is similar to the hexagonal lattice.
\end{lemma}

This is the value $\gamma_2=2/\sqrt3$ of Hermite's constant in dimension two; see~\cite[Chapter~II]{Cas59}. It follows
from reduction of binary quadratic forms: a reduced form $ax^2+bxy+cy^2$ with $|b|\le a\le c$ satisfies
$a^2\le ac\le\frac43(ac-b^2/4)$, where $a$ is the squared length of a shortest nonzero vector and $ac-b^2/4=d^2$.

\begin{proof}[Proof of Theorem~\ref{thm:main}]
A nonzero vector of $\Lambda$ with $m=0$ is $-ju_2$ with $j\ne0$, and its squared length is $4j^2/A\ge4/A>4$. By
Lemma~\ref{lem:hermite} applied with $d=2$, $\Lambda$ has a nonzero vector of squared length at most $4/\sqrt3<4$, and
this vector therefore has $m\ne0$. Hence $\mu\le4/\sqrt3$, and~\eqref{eq:lattice} gives
\[
\gamma(n,k)\le\frac{\pi^2\sqrt k}{2n^2}\cdot\frac{4}{\sqrt3}=\frac{2\pi^2}{\sqrt3}\,\frac{\sqrt k}{n^2}.
\]
The statement about $\gap(n,k)$ follows from Theorem~\ref{thm:apw}, since its error factor $1+o(1)$ is uniform in $k$.
\end{proof}

\begin{remark}[Equality]\label{rem:equality}
Take $n=26$ and $k=12$, so that $\theta=7/13$ and $A=\sqrt{12}/26=\sqrt3/13$. For $m=3$ we have $3\theta=21/13$ and
$b_3=5/13$, so
\[
A\cdot 9+\frac{b_3^2}{A}=\frac{9\sqrt3}{13}+\frac{25}{13\sqrt3}=\frac{52}{13\sqrt3}=\frac{4}{\sqrt3}.
\]
Hence $\gamma(n,k,3)$ equals the right side of~\eqref{eq:bound}, and since $\gamma(n,k)$ cannot exceed it,
$\gamma(26,12)=2\pi^2\sqrt{12}/(\sqrt3\cdot26^2)=\pi^2/169$. So~\eqref{eq:bound} is attained at a finite $(n,k)$, as well as
asymptotically along the sequence $k=n-(2/3^{1/4}+o(1))n^{1/4}$ noted in~\cite[Section~6]{APW08}.
Theorem~\ref{thm:main} gives only the upper bound for $\gap(n,k)$; that the envelope
is approached along $k=(\alpha+o(1))n$ for every $\alpha\in(0,1]$, as indicated in~\cite[Section~6]{APW08}, is not
addressed here.
\end{remark}

\section{Numerical checks}\label{sec:checks}

The program in the repository~\cite{Repo} evaluates $\gamma(n,k)$ from~\eqref{eq:gamma} and the exact gap from the
$n\times n$ transition matrix of the single-card chain. Over all $2\le n\le2000$ and $1\le k<n$, the largest value of
$\gamma(n,k)\,n^2/\sqrt k$ is $11.396437516$, which agrees with $2\pi^2/\sqrt3=11.3964375155\dots$ to the printed
digits and occurs at $(n,k)=(26,12)$. For $n=1500$ and $k=1,38,75,\dots$, the largest value of
$\gap(n,k)\,n^2/\sqrt k$ is $11.106$; for $n=300$ and $k=1,24,47,\dots$ it is $8.644$. At $k=1$ the ratio
$\gap(n,k)/\gamma(n,k)$ is $0.9995$ for $n=300$. Over the values of $k$ checked for $n=1500$ the ratio lies between
$0.79$ and $1.025$ (the largest at $k=260$). The small ratios occur when the minimising $m$ in~\eqref{eq:gamma} has
$b_m$ not small, where the quadratic expression is a poor approximation; in the cases computed they move toward $1$
as $n$ grows. For $k/n=0.79$
the ratio is $0.794$, $0.791$, $0.991$ and $0.975$ at $n=750$, $1500$, $3000$ and $6000$, and for $k/n=0.7407$ it is
$0.741$, $1.005$ and $1.002$ at $n=750$, $1500$ and $3000$.

\section*{Acknowledgements and funding}
This work was done by multiple AI agents, Claude (Anthropic) and Codex (OpenAI), connected to
Math Brain~\cite{MathBrainRepo} and working together inside the Brain framework~\cite{BrainFramework}, a shared,
persistent memory in which the agents record claims,
proofs, failed attempts and methods, retrieve one another's results, coordinate who works on what, and audit each
other's work. For more on the Brain and Math Brain, see the Math Brain paper~\cite{MathBrainPaper}, the Math Brain
repository~\cite{MathBrainRepo} and the Brain framework repository~\cite{BrainFramework}.

The research presented in this paper was supported by the European Research Council (ERC) under the European Union's
Horizon Europe research and innovation programme (grant agreement No.~101041711), by the Simons Foundation as part of
the Collaboration on the Mathematical and Scientific Foundations of Deep Learning, by Heights Labs, by Convex Nexus
Capital, by the Israel Science Foundation (grant number 2258/19), and by the Israel Science Foundation (ISF Grant
4101/25). Computational resources were provided in part by Google Cloud through its education credits programme
(credit 530550103).

\section*{Declaration of competing interest}
The author holds equity in Heights Labs and in Convex Nexus Capital and has a management role at Convex Nexus Capital.
These entities are included in the funding declaration above.

\section*{Data availability}
The program described in Section~\ref{sec:checks} is in the repository~\cite{Repo}.

\section*{Declaration of generative AI and AI-assisted technologies}
This paper is part of a broader research project on AI proving problems. Claude agents (Anthropic) proved
Theorem~\ref{thm:main}, wrote the program of Section~\ref{sec:checks} and prepared the manuscript. The shuffle's
single-card analysis, the quantity $\gamma(n,k)$, Theorem~\ref{thm:apw}, the bound
$\gamma(n,k)\le2\pi^2k^{1/2}n^{-2}$ from Dirichlet's theorem, the bound $(2\pi^2+o(1))n^{-3/2}$, the asymptotic
sharpness of the constant $2\pi^2/\sqrt3$ for $\gamma$ and the question answered here are due to Angel, Peres and
Wilson~\cite{APW08}. Lemma~\ref{lem:hermite} is classical. The new contribution is the identification of $\gamma(n,k)$ with the minimum of a
plane lattice of determinant $2$, which improves the constant to $2\pi^2/\sqrt3$ and gives Theorem~\ref{thm:main}. The author directed the work and is responsible for the paper.

\begin{thebibliography}{99}
\bibitem{APW08} O.~Angel, Y.~Peres and D.~B.~Wilson, \emph{Card shuffling and Diophantine approximation}, Ann. Appl.
Probab. \textbf{18} (2008), 1215--1231, \doi{10.1214/07-AAP484}, \arxiv{0707.2994}.
\bibitem{BMO23} O.~Blumberg, B.~Morris and H.~Oberschelp, \emph{Mixing time of the overlapping cycles shuffle}, 2023,
\arxiv{2310.15438}.
\bibitem{Cas59} J.~W.~S.~Cassels, \emph{An Introduction to the Geometry of Numbers}, Grundlehren der mathematischen
Wissenschaften~99, Springer, Berlin, 1959.
\bibitem{Jon06} J.~Jonasson, \emph{Biased random-to-top shuffling}, Ann. Appl. Probab. \textbf{16} (2006), 1034--1058, \doi{10.1214/10505160600000097}.
\bibitem{Jon11} J.~Jonasson, \emph{Mixing time bounds for overlapping cycles shuffles}, Electron. J. Probab.
\textbf{16} (2011), 1281--1295, \doi{10.1214/EJP.v16-912}.
\bibitem{Repo} Y.~Shmalo, \emph{Overlapping-cycles single-card gap bound: paper and check program}, 2026,
release v1.0, \url{https://github.com/yspennstate/overlapping-cycles-gap-bound/tree/v1.0}.
\bibitem{MathBrainRepo} Y.~Shmalo, \emph{Math Brain: public snapshot, code and tests}, 2026,
\url{https://github.com/yspennstate/math-brain-public/tree/8c9a6bd478c4ec4082347151b7ae16db50c917dd}.
\bibitem{MathBrainPaper} Y.~Shmalo, \emph{Math Brain: shared, auditable memory for multi-agent mathematical research},
October 2026, \doi{10.5281/zenodo.24150837},
\url{https://github.com/yspennstate/math-brain-public/tree/main/papers/math-brain-2026-10-24}.
\bibitem{BrainFramework} Y.~Shmalo, \emph{Brain Framework}, public software repository, 2026,
\url{https://github.com/yspennstate/brain-framework-public/tree/5629e644cdd510173ae97e9aa92642dabd22dc57}.
\end{thebibliography}

\end{document}
